% Fonte principal: abstract_pt e keywords_pt em tcc.conf.
\begin{resumo}
O mapeamento geológico do Brasil ainda é incompleto, e a prospecção mineral tradicional depende de levantamentos de campo e análises laboratoriais demorados e custosos. Dados espectrais de satélite permitem observar grandes áreas e podem apoiar a seleção de locais para investigação. Este trabalho propõe avaliar o TabPFN, um \textit{foundation model} para dados tabulares pré-treinado com conjuntos sintéticos, na previsão de ocorrências minerais e na predição da composição química a partir de dados espectrais de satélite. Para isso, será construído um conjunto de dados que associa a substância registrada nos processos minerários ativos da Agência Nacional de Mineração (ANM) e as referências geológicas e geoquímicas do Serviço Geológico do Brasil (SGB) aos valores de reflectância de imagens Sentinel-2 de nível 2A. Na previsão de ocorrências, um único modelo de classificação atribuirá probabilidades às substâncias; na predição da composição química, condicionada à disponibilidade de referências adequadas, será ajustado um modelo de regressão para cada elemento. O TabPFN será comparado a modelos utilizados nos trabalhos relacionados, como \textit{Random Forest} e XGBoost, por meio de validação cruzada por blocos espaciais em todo o Brasil e de métricas de classificação e de regressão. Um estudo de ablação verificará a contribuição de cada grupo de atributos. Os conjuntos derivados cuja redistribuição for permitida serão disponibilizados no Kaggle. Espera-se verificar se a abordagem pode apoiar a triagem inicial de áreas para investigação, sem substituir a confirmação geológica em campo.
\end{resumo}
